\scanpage{27}
\mspage{\textcircled{13}}
De façon précise, \inserted{la donnée} des \inserted{$M_0$-torseurs} $M_i$ et \inserted{de} l'isom. can. $6$ revient à la donnée du torseur $P_M$ sous $V_J(M_0)$, avec
\[
(9)\qquad M_i\simeq P_M\wedge_{M_0^J}(M_0,\operatorname{pr}_i)
\qquad\text{pour }i\in J
\]
où $\operatorname{pr}_i:M_0^J\mathbin{-}M_0$ la projection d'indice $i\in J$.

Notons qu'on a des iso. canoniques (valables pour $\forall$ vectoriel de rang $2$ sur $\mathbb{F}_2$)
\[
(10)\qquad
\operatorname{End}(M_0)\simeq V_J(M_0)\xrightarrow{\sim}M_0\otimes M_0
\simeq M_0\oplus\mathbb{F}_2^{\underline{\omega}_{M_0}}
\]
\[
\underline{\omega}_{M_0}\simeq\underline{\omega}
(\underbrace{M_0^*}_{J})\quad\text{ens des orientations de }J.
\]
\margin{\uline{NB} Pour $\forall$ vectoriel $W$ sur $\mathbb{F}_2$ on a
\[
(11)\qquad\operatorname{Hom}_{\mathbb{F}_2}(M_0,W)\simeq V_J(W)
\]
Car $M_0\simeq\mathbb{F}_2^J/\mathbb{F}_2\text{ diag}$}
\diagram{\begin{tikzcd}
(12)\quad\operatorname{End}(M_0) \arrow[r,"\sim"] & V_J(M_0)\\
u \arrow[r,mapsto] &(u(x_i))_{i\in J}
\end{tikzcd}}
\[
\operatorname{End}(M_0)\simeq M_0^\vee\otimes M_0\simeq M_0\otimes M_0
\quad\text{car }M_0^\vee\simeq M_0
\]
(valable pour $\forall$ vect. de rang $2$ sur $\mathbb{F}_2$)

\diagram{\begin{tikzcd}
(13)\quad M_0 \arrow[r,hook] & \operatorname{End}(M_0)\\
M_0^* \arrow[u,phantom,"\cup" description] \arrow[r,hook] &
\operatorname{Aut}(M_0)\simeq\unsure{\mathcal{G}_J|\mathfrak{S}_J}
\arrow[u,phantom,"\cup" description]
\end{tikzcd}}
\[
x_i\longmapsto\deleted{\text{transposition}}\inserted{\sigma_i}
\quad\text{réflexion de $M_0$ qui fixe $x_i$}
\]
(donc \deleted{\unsure{qui est}}\inserted{transposition} des $x_j$ et $x_k$)
[\uline{NB} Comme $\sum_{i\in J}\sigma_i=0$, cela correspond à hom
$M_0\longrightarrow\operatorname{End}(M_0)$]

\diagram{\begin{tikzcd}
(14)\quad\mathbb{F}_2^{\underline{\omega}_{M_0}} \arrow[r,hook] &
\operatorname{End}(M_0)\\
\underline{\omega}_{M_0} \arrow[u,phantom,"\cup" description] \arrow[r,hook,"\mathrm{can}"] &
\operatorname{Aut}(M_0)\simeq\unsure{\mathcal{G}_J|\mathfrak{S}_J}
\arrow[u,phantom,"\cup" description]\\
\underline{\omega}_J
\simeq\unsure{\mathcal{G}_J^+|\mathfrak{S}_J^+}\setminus\{\operatorname{id}_J\}
\arrow[u,phantom,"\Vert" description] &
\end{tikzcd}}
\margin{\uline{NB} \uline{On a} -- si $\rho\in\underline{\omega}$,
$\rho'\underset{\mathrm{def}}{=}\rho^{-1}$
\[
\rho+\rho'+\unsure{1}\operatorname{id}=0
\]
i.e. $\rho+\rho'=\operatorname{id}_{M_0}$.}
i.e. l'hom $\mathbb{F}_2^{\underline{\omega}}\hookrightarrow\operatorname{End}(M_0)$ applique $1$ sur $1$.

Si $i,j\in J$, $i\ne j$, alors $\sigma_i,\sigma_j,\rho,\rho'\in\operatorname{End}(M_0)$ sont lin. indépendants sur $\mathbb{F}_2$ donc forment une base, \deleted{\unsure{de}} dès que $x_i,x_j$ en forment une de $M_0$, et $\rho,\rho'$ de $\mathbb{F}_2^{\underline{\omega}}$\ldots
% archive foliation: 26
